\chapter*{Abbreviations and Notation}
\addcontentsline{toc}{chapter}{Abbreviations and Notation}

This chapter defines the abbreviations and mathematical notation used throughout the report. The notation table is intended to be exhaustive for the current version. When new symbols are introduced in later experiments, they should be added here and also explained at their first use in the corresponding chapter.

\section*{Abbreviations}
\begin{longtable}{>{\bfseries}p{0.20\textwidth}p{0.72\textwidth}}
\toprule
Abbreviation & Meaning\\
\midrule
\endfirsthead
\toprule
Abbreviation & Meaning\\
\midrule
\endhead
AACOPF & Adaptive Ant Colony Optimization Particle Filter\\
ACO & Ant Colony Optimization\\
DR & Dead Reckoning\\
DW3000 & Qorvo DW3000 ultra-wideband transceiver family used in the intended hardware platform\\
ESP32 & Espressif ESP32 microcontroller family used in the intended embedded platform\\
IMU & Inertial Measurement Unit\\
INS & Inertial Navigation System\\
LOS & Line of Sight\\
NLOS & Non-Line of Sight\\
PF & Particle Filter\\
p95 & 95th percentile of an error distribution\\
RMSE & Root-Mean-Square Error\\
RW & Random Walk\\
UWB & Ultra-Wideband\\
YAML & YAML Ain't Markup Language; text format used for version-controlled experiment configuration\\
\bottomrule
\end{longtable}

\section*{Notation conventions}
\begin{longtable}{p{0.24\textwidth}p{0.68\textwidth}}
\toprule
Symbol / convention & Meaning\\
\midrule
\endfirsthead
\toprule
Symbol / convention & Meaning\\
\midrule
\endhead
$k$ & Discrete-time sample index.\\
$i$ & Particle index in the particle filter.\\
$t$ & Continuous time in seconds.\\
$\Delta t$ & Sampling interval between discrete states.\\
$n$ & Navigation frame: fixed two-dimensional world frame in which global position is expressed.\\
$b$ & Body frame: coordinate frame rigidly attached to the mobile node and IMU.\\
$(\cdot)^n$ & Quantity expressed in the navigation frame.\\
$(\cdot)^b$ & Quantity expressed in the body frame.\\
$(\cdot)_m$ & Measured sensor quantity.\\
$\hat{(\cdot)}$ & Estimated quantity.\\
$\bm{(\cdot)}$ & Vector or matrix quantity.\\
$(\cdot)^\top$ & Matrix/vector transpose.\\
$\lVert\cdot\rVert_2$ & Euclidean norm.\\
$\mathcal{N}(\mu,\sigma^2)$ & Gaussian distribution with mean $\mu$ and variance $\sigma^2$. For vector quantities, a covariance matrix is given explicitly when needed.\\
$\delta(\cdot)$ & Dirac delta distribution used in the particle representation of a probability density.\\
$\operatorname{wrap}(\theta)$ & Mapping of angle $\theta$ to the interval $[-\pi,\pi)$.\\
\bottomrule
\end{longtable}

\section*{State, motion, and IMU quantities}
\begin{longtable}{p{0.30\textwidth}p{0.62\textwidth}}
\toprule
Symbol & Meaning\\
\midrule
\endfirsthead
\toprule
Symbol & Meaning\\
\midrule
\endhead
$\bm{x}_k$ & Navigation state at sample $k$, currently $[p_{x,k},p_{y,k},v_{x,k},v_{y,k},\psi_k]^\top$.\\
$p_{x,k}$ & Global $x$-position of the target node in the navigation frame at sample $k$.\\
$p_{y,k}$ & Global $y$-position of the target node in the navigation frame at sample $k$.\\
$\bm{p}_k$ & Target position vector $[p_{x,k},p_{y,k}]^\top$.\\
$v_{x,k}$ & Target velocity component along the navigation-frame $x$-axis.\\
$v_{y,k}$ & Target velocity component along the navigation-frame $y$-axis.\\
$\bm{v}_k$ & Navigation-frame velocity vector $[v_{x,k},v_{y,k}]^\top$.\\
$\psi_k$ & Target yaw angle at sample $k$; orientation of the body frame relative to the navigation frame.\\
$R_b^n(\psi)$ & Two-dimensional rotation matrix that maps a body-frame vector into the navigation frame for yaw angle $\psi$.\\
$\bm{u}_k$ & Ideal planar IMU input vector $[a_{x,k}^b,a_{y,k}^b,\omega_{z,k}]^\top$.\\
$a_{x,k}^b$ & Ideal body-frame longitudinal acceleration at sample $k$.\\
$a_{y,k}^b$ & Ideal body-frame lateral acceleration at sample $k$.\\
$\bm{a}_k^b$ & Ideal body-frame horizontal acceleration vector $[a_{x,k}^b,a_{y,k}^b]^\top$.\\
$\omega_{z,k}$ & Ideal angular rate about the vertical/body $z$-axis, i.e. yaw rate.\\
$\bm{a}_{m,k}^b$ & Measured body-frame horizontal acceleration vector.\\
$\omega_{m,k}$ & Measured yaw rate.\\
$\bm{b}_{a,k}$ & Additive accelerometer-bias vector at sample $k$.\\
$b_{g,k}$ & Additive gyroscope yaw-rate bias at sample $k$.\\
$\bm{n}_{a,k}$ & Accelerometer white-noise sample.\\
$n_{g,k}$ & Gyroscope white-noise sample.\\
$\bm{w}_{ba,k}$ & Accelerometer-bias random-walk increment between samples $k$ and $k+1$.\\
$w_{bg,k}$ & Gyroscope-bias random-walk increment between samples $k$ and $k+1$.\\
$\psi_{k+1/2}$ & Midpoint yaw used to rotate acceleration during the interval $[k,k+1]$.\\
$\bm{a}_k^n$ & Measured horizontal acceleration after rotation into the navigation frame.\\
$b_a$ & Generic constant acceleration bias used only in the drift-growth illustration.\\
$\delta p(t)$ & Approximate position error induced by a constant acceleration bias after elapsed time $t$.\\
$s(t)$ & Magnitude of the navigation-frame velocity in the deterministic Phase-1 trajectory, $s(t)=\lVert\bm{v}(t)\rVert_2$.\\
\bottomrule
\end{longtable}

\section*{UWB quantities}
\begin{longtable}{p{0.30\textwidth}p{0.62\textwidth}}
\toprule
Symbol & Meaning\\
\midrule
\endfirsthead
\toprule
Symbol & Meaning\\
\midrule
\endhead
$\bm{p}_{A,k}$ & Known navigation-frame position of auxiliary UWB node $A$ at sample $k$.\\
$\bm{p}_{A}(t)$ & Continuous-time trajectory used to specify auxiliary-node motion in a simulation scenario.\\
$d_k$ & Ideal geometric range between target and auxiliary node at sample $k$.\\
$z_k$ & Measured UWB range at sample $k$.\\
$\nu_k$ & Additive UWB range-measurement noise.\\
$\sigma_{\mathrm{UWB}}$ & Standard deviation of the Gaussian UWB range-noise model.\\
\bottomrule
\end{longtable}

\section*{Particle-filter quantities}
\begin{longtable}{p{0.30\textwidth}p{0.62\textwidth}}
\toprule
Symbol & Meaning\\
\midrule
\endfirsthead
\toprule
Symbol & Meaning\\
\midrule
\endhead
$N_p$ & Number of particles.\\
$\bm{x}_k^{(i)}$ & State of particle $i$ at sample $k$.\\
$\bm{p}_k^{(i)}$ & Position component of particle $i$.\\
$w_k^{(i)}$ & Normalized importance weight of particle $i$.\\
$\bm{\epsilon}_{a,k}^{(i)}$ & Additional sampled acceleration perturbation used to propagate particle $i$.\\
$\epsilon_{\omega,k}^{(i)}$ & Additional sampled yaw-rate perturbation used to propagate particle $i$.\\
$\sigma_{a,\mathrm{PF}}$ & Standard deviation of each component of the particle acceleration perturbation.\\
$\sigma_{\omega,\mathrm{PF}}$ & Standard deviation of the particle yaw-rate perturbation.\\
$\hat d_k^{(i)}$ & UWB range predicted from particle $i$.\\
$r_k^{(i)}$ & UWB residual $z_k-\hat d_k^{(i)}$ for particle $i$.\\
$N_{\mathrm{eff},k}$ & Effective sample size of the particle weights at sample $k$.\\
$\gamma$ & Resampling fraction; resampling is triggered when $N_{\mathrm{eff},k}<\gamma N_p$.\\
$\bm{\sigma}_0$ & Vector of standard deviations used for the Gaussian initial particle cloud in the known-pose baseline.\\
$\hat{\bm{x}}_k$ & Particle-filter state estimate.\\
$\hat{\bm{p}}_k$ & Estimated position component of $\hat{\bm{x}}_k$.\\
$\hat{\psi}_k$ & Estimated yaw angle, computed using a weighted circular mean.\\
\bottomrule
\end{longtable}

\section*{Observability and evaluation quantities}
\begin{longtable}{p{0.30\textwidth}p{0.62\textwidth}}
\toprule
Symbol & Meaning\\
\midrule
\endfirsthead
\toprule
Symbol & Meaning\\
\midrule
\endhead
$F_k$ & Local state-transition Jacobian from sample $k$ to $k+1$.\\
$H_k$ & Local Jacobian of the scalar UWB range measurement with respect to the state.\\
$I_q$ & $q\times q$ identity matrix.\\
$0_{r\times c}$ & $r\times c$ zero matrix.\\
$\bm{a}_{\psi,k}$ & Derivative of navigation-frame acceleration with respect to yaw.\\
$\Phi_{j,0}$ & State-transition sensitivity from sample $0$ to sample $j$, equal to $F_{j-1}\cdots F_0$ for $j>0$.\\
$\mathcal{O}_{0:K}$ & Finite-horizon local observability matrix from sample $0$ through $K$.\\
$\sigma_{\min}(\mathcal{O})$ & Smallest singular value of the observability matrix.\\
$\sigma_{\max}(\mathcal{O})$ & Largest singular value of the observability matrix.\\
$e_{p,k}$ & Euclidean position error $\lVert\hat{\bm{p}}_k-\bm{p}_k\rVert_2$.\\
$e_{\psi,k}$ & Wrapped yaw error $\operatorname{wrap}(\hat\psi_k-\psi_k)$.\\
$N_s$ & Number of samples included in an evaluation metric.\\
$\mathrm{RMSE}_p$ & Root-mean-square position error.\\
$\mathrm{RMSE}_{\psi}$ & Root-mean-square yaw error.\\
$Q_{0.95}$ & Empirical 95th-percentile/quantile operator.\\
\bottomrule
\end{longtable}

\section*{Phase-1 noise and reference-paper quantities}
\begin{longtable}{p{0.30\textwidth}p{0.62\textwidth}}
\toprule
Symbol & Meaning\\
\midrule
\endfirsthead
\toprule
Symbol & Meaning\\
\midrule
\endhead
$\sigma_a$ & Per-sample standard deviation of simulated accelerometer white noise.\\
$\sigma_{\omega}$ & Per-sample standard deviation of simulated gyroscope white noise.\\
$\sigma_{ba}$ & Accelerometer-bias random-walk intensity; the bias increment standard deviation over one step is $\sigma_{ba}\sqrt{\Delta t}$.\\
$\sigma_{bg}$ & Gyroscope-bias random-walk intensity; the bias increment standard deviation over one step is $\sigma_{bg}\sqrt{\Delta t}$.\\
$\Delta L$ & Travelled-distance increment supplied by the higher-level DR system in the Han et al. reference formulation.\\
$\Delta\psi$ & Azimuth/yaw increment supplied by the higher-level DR system in the Han et al. reference formulation.\\
\bottomrule
\end{longtable}
